\chapter{Conclusiones}\label{sec.ej1:conclusiones}

La implementación y verificación de la unidad de procesamiento con arquitectura Harvard desarrollada en este trabajo
práctico dio resultados coherentes con los requerimientos teóricos planteados para el diseño de procesadores.

En primer término, la organización jerárquica del camino de datos permitió desacoplar con éxito la memoria interna
(banco de registros de tres puertos) de las unidades aritméticas y lógicas de procesamiento. La utilización de puertos
de lectura puramente combinacionales y un único puerto de escritura síncrono facilitó la realización de ciclos de
cálculo en los cuales los operandos se leen, se operan en la ALU o en el desplazador, y el resultado queda disponible
en el bus de escritura dentro del mismo período de reloj, efectuándose la captura en memoria al presentarse el flanco
activo.

En segundo lugar, se validó la eficacia de la descomposición de la palabra de control (\texttt{regCTRL}). Esta práctica
resultó indispensable para verificar visualmente en el esquemático RTL cómo la palabra de control ingresa y bifurca sus
campos hacia las instancias de procesamiento y los multiplexores de selección.

En tercer lugar, la simulación temporal en el banco de pruebas confirmó la ejecución exacta de las operaciones
solicitadas:
\begin{itemize}
  \item La adición entre los valores precargados en $R_0$ (5) y $R_1$ (3) produjo el resultado 8 y
        activó de manera correcta la bandera de signo negativo ($N = 1$), debido a la convención de signo en el bit 3.
  \item La sustracción entre el acumulador $R_2$ (8) y el dato constante externo $cIN$ (2) a través de la vía de
        entrada alternativa del multiplexor B entregó el resultado 6 (\texttt{0110b}) con activación de la bandera de
        acarreo ($C = 1$), ratificando la adecuada implementación del complemento a dos ($A + \sim B + 1$).
  \item La posterior operación de desplazamiento lógico a la izquierda duplicó el valor almacenado en $R_2$ alcanzando
        12 (\texttt{1100b}), verificando la correcta selectividad del multiplexor de la unidad funcional.
\end{itemize}

En síntesis, la experiencia permitió comprobar de manera analítica y experimental el funcionamiento coordinado de una
unidad de procesamiento Harvard, sentando las bases de hardware sobre las cuales se diseñan las unidades de control y
los secuenciadores de instrucciones en arquitecturas computacionales.
